% Worked Examples: Concept 1.1.1 - Introduction
\documentclass[11pt,a4paper]{article}
\usepackage{worked-examples-template}

\title{\textbf{Worked Examples}\\[0.5em]
\large Concept 1.1.1: Introduction\\[0.3em]
\normalsize \coursesubtitle{} -- \coursetitle}
\author{Assoc.\ Prof.\ William Robertson \and Dr Sean McGowan}
\date{\today}

\begin{document}

\maketitle
\tableofcontents
\newpage

\section{Concept 1.1.1: Introduction}

\subsection{Example 1: Cruise Control System - Open-Loop vs Closed-Loop}

\paragraph{Problem:} A car's cruise control system aims to maintain a constant speed of $v_{\text{ref}} = 100$ km/h on a highway. Compare two approaches:
\begin{itemize}
\item \textbf{Open-loop control:} The throttle is set to a fixed position that achieves 100 km/h on flat ground.
\item \textbf{Closed-loop control:} The throttle adjusts continuously based on the error between the desired and actual speed.
\end{itemize}

Analyze what happens when the car encounters an uphill slope (5\% grade) that creates an additional resistance force of 500 N. The car has mass $m = 1500$ kg and aerodynamic drag coefficient $c = 30$ N·s/m.

\paragraph{Solution:}

\textbf{Open-Loop Control:}

In open-loop control, the throttle setting $u_0$ is fixed. At the desired speed on flat ground:
\begin{align}
F_{\text{drive}} = F_{\text{drag}} = cv_{\text{ref}} = 30 \times \frac{100}{3.6} = 833 \text{ N}
\end{align}

When encountering the uphill slope, the total resistance becomes:
\begin{align}
F_{\text{resistance}} = 833 + 500 = 1333 \text{ N}
\end{align}

Since the drive force remains constant at 833 N, there is a net deceleration:
\begin{align}
ma = 833 - 1333 = -500 \text{ N} \quad \Rightarrow \quad a = \frac{-500}{1500} = -0.33 \text{ m/s}^2
\end{align}

The car will slow down continuously until reaching a new equilibrium where:
\begin{align}
F_{\text{drive}} = cv_{\text{new}} + 500 \quad \Rightarrow \quad v_{\text{new}} = \frac{833 - 500}{30} = 11.1 \text{ m/s} = 40 \text{ km/h}
\end{align}

\textbf{Result:} The speed drops from 100 km/h to 40 km/h — a significant error with no automatic correction.

\textbf{Closed-Loop Control:}

In closed-loop control, a controller measures the speed error $e = v_{\text{ref}} - v$ and adjusts the throttle. Using a simple proportional-integral (PI) controller:
\begin{align}
u(t) = K_p e(t) + K_i \int_0^t e(\tau) d\tau
\end{align}

When the car starts slowing down due to the slope:
\begin{enumerate}
\item The error $e = 100 - v$ becomes positive
\item The controller increases throttle to provide more drive force
\item The integral term accumulates, ensuring zero steady-state error
\item The car returns to and maintains 100 km/h
\end{enumerate}

With appropriate controller gains (e.g., $K_p = 50$, $K_i = 20$), the system can reject the disturbance and maintain the reference speed within 1-2\% error.

\textbf{Key Insights:}
\begin{itemize}
\item \textbf{Open-loop systems} are simple but cannot compensate for disturbances or model uncertainties
\item \textbf{Closed-loop (feedback) systems} continuously measure the output and correct for errors
\item Feedback is essential when the system is subject to unknown disturbances or parameter variations
\item The cost of feedback is increased complexity (sensors, controllers) and potential stability issues if not designed properly
\end{itemize}

This example illustrates the fundamental motivation for feedback control systems in real-world applications.
\end{document}
